\documentclass[12pt]{article}
\usepackage[utf8]{inputenc}
\usepackage{amsmath,amssymb,amsthm}
\usepackage{geometry}
\geometry{a4paper, margin=1in}

\title{Project Euler Problem 566: Cake Icing Puzzle}
\author{}
\date{}

\begin{document}
\maketitle

\section{Problem Statement}
Adam plays a game with his birthday cake. He cuts circular sectors of different angles and flips them upside down. Three cutting angles are used cyclically:
\[
x = \frac{360}{a}, \quad y = \frac{360}{b}, \quad z = \frac{360}{\sqrt{c}}
\]
Define $F(a,b,c)$ as the minimum number of flips to restore all icing to the top. Define:
\[
G(n) = \sum_{\substack{9 \le a < b < c \le n}} F(a,b,c)
\]
Find $G(53)$.

\section{Mathematical Analysis}

\subsection{State Representation}
Model the cake as the circle $[0, 360)$. Each point has state $+1$ (icing up) or $-1$ (icing down). A flip of sector $[\theta_0, \theta_0 + \alpha)$ negates states in that arc, followed by rotation by $\alpha$.

\subsection{Simulation Approach}
The state is tracked by a sorted list of breakpoints on the circle. Between consecutive breakpoints, the icing state is uniform. Each flip operation:
\begin{enumerate}
    \item Introduces at most 2 new breakpoints (at the sector boundaries)
    \item Toggles the state of intervals within the sector
    \item Rotates all breakpoints by the angle
\end{enumerate}

\subsection{Termination}
The process terminates when all intervals have state $+1$. For the combination of rational and irrational angles, the Poincar\'e recurrence theorem guarantees eventual return to the initial state.

\section{Algorithm}
For each triple $(a,b,c)$:
\begin{enumerate}
    \item Initialize circle with uniform state $+1$
    \item Cyclically apply flips with angles $360/a$, $360/b$, $360/\sqrt{c}$
    \item Track state changes using interval breakpoints with high-precision arithmetic
    \item Count total flips until restoration
\end{enumerate}

\section{Complexity}
$O\!\left(\binom{n-8}{3} \cdot L\right)$ where $L$ is the average number of flips per triple.

\section{Answer}

\[
\boxed{329569369413585}
\]

\end{document}
